% !TeX program = xelatex
% !BIB program = biber
\documentclass[12pt,letterpaper]{article}
\usepackage{fontspec}
\setmainfont{Liberation Serif} % Metric compatible with Times New Roman
\usepackage[margin=1in]{geometry}
\usepackage{setspace}
\doublespacing 
\usepackage{indentfirst} 
\usepackage[american]{babel}
\usepackage{csquotes}
\usepackage[style=mla, backend=biber]{biblatex}
\addbibresource{bib.bib} % Bibliography file
\usepackage{fancyhdr}
\pagestyle{fancy}
\fancyhf{} 
\fancyhead[R]{Last \thepage}
\renewcommand{\headrulewidth}{0pt} % Remove header line
\setlength{\bibhang}{0.5in} % Make Works Cited use half-inch indents
\renewenvironment{quote}%
    {\list{}{\leftmargin=0.5in \rightmargin=0in}\item\relax}{\endlist}
 % Packages and settings
\fancyhead[R]{Last \thepage}

\begin{document}

\noindent First Last\\
\noindent Professor\\
\noindent Aesthetics and Imagination\\
\noindent 19 April 2026

% Center environment adds unwanted vertical padding
{\centering Art, Cognition, and Entangled Domains\par}

% --- Content ---

In \textit{The Entanglement}, Alva Noë applies and develops the concept of
``entanglement''; concluding, at one point, that ``Art and philosophy are
entangled with life and science'' \autocite[202]{noeEntanglement2023}. Here, I
argue that this entanglement is broader than the particular fields Noë mentions;
individual disciplines or concepts are not the only kinds of things subject to
the entanglement. Instead, entire domains or spheres of human existence are
entangled. In particular, using the work of Hannah Arendt and Byung-Chul Han, I
develop the concept of an exclusive sphere that stands outside the public sphere
of common human affairs. These two spheres turn out to be entangled; moreover,
this political entanglement can explain when and how the exclusive becomes
politically salient.

In my analysis, I focus on art, cognitive science, and the cognitive science of
art. How each item relates to the political realm (that is, whether it respects
its entanglement) predicts its political power. At times, my interpretation of
the entanglement differs from Noë's. Regardless, the points we share are
fundamental to the entanglement. These points are exemplified by the
entanglement of dance (a common, ``habitual'' activity, in Noë's terms) and
Dance (an artistic practice related to choreography).

Lowercase ``dance,'' as Noë describes it, is an ostensibly naïve and
unencumbered activity; dancing is doing what feels right in the moment.
Moreover, dance is habitual in the sense that, to ``the degree to which we
excel, we do so by letting ourselves go'' \autocite[27]{noeEntanglement2023}.
Noë's concept of habit is slightly idiosyncratic; it concerns an automatic way
of being that is more susceptible to inclination. This authentic, habitual way
of being is what Noë calls the first-order self
\autocite[19]{noeEntanglement2023}.  By contrast, uppercase ``Dance'' (or
choreography) involves our second-order selves. Dance requires reflection on
dance so that it may put the latter on display and investigate it
\autocite[29]{noeEntanglement2023}. Thus, Dance requires us to distance
ourselves from the initially familiar act of dancing.

It is worth noting the role of distance in the relationship between dance and
Dance. In \textit{The Entanglement}, this concept is hardly central, but I will
later argue that it is essential for a politically salient entanglement. For
Han, ``Respect presupposes a distanced look''; to respect something is, in part,
to distance yourself from it. Han contrasts respect with ``spectacle,'' a kind
of ``voyeuristic gazing'' without deference \autocite[1]{hanSwarm2017}. When one
engages in Dance, they remove themselves from the immediacy of dance while
keeping it as their focus. This sort of healthy or reserved looking, I argue, is
essential to the broader entanglements discussed down the line.

% point out two things: first and second order are bi-directional and first not
% naive also, does not eliminate differences between entangled concepts
Returning to dance and Dance, it is crucial to recognize that the relationship
between first- and second-order acts is bi-directional. While the first-order
gives the second-order its substance, the second-order also alters the first.
When one dances, they do so intuitively; nonetheless, dancing calls on
choreography. Individual dance moves, as well as the norms governing what
constitutes dance, will, in part, come from choreography or Dance
\autocite[30]{noeEntanglement2023}. In short, the first-order act of dancing was
not naïve or fully natural. Regardless, the entanglement does not collapse the
conceptual distinction between dance and choreography; in fact, Noë argues that
this conceptual distinction is a condition of their entanglement
\autocite[31]{noeEntanglement2023}.

Dance and dance are cultural practices, so they make for good examples of an
entanglement. By their nature, cultural practices will involve standards
(however fuzzy or implicit) that help us determine what actions do and do not
count as belonging to the practice. When we find that two practices are
relevantly similar, we find it natural to extend the standards of one area to
another. In the case of dance and Dance, the latter's technicality and
artfulness are relevant enough to merit a new category. Nevertheless, we
maintain that dance and Dance are concerned with the same kind of thing, namely,
expressive movement. Both of these views exist in a common understanding of
dancing, but the entanglement makes their relationship explicit.

As is popular to assert nowadays, one could argue that the entanglement of dance
and Dance is best understood in terms of a spectrum. It could be, for instance,
that Dance is dance made technical. Consequently, the entanglement is
unnecessary to understand how these concepts relate to each other; instead, we
only need a spectrum of artfulness or technicality. That dance and Dance exist
along the same spectrum may capture some aspects of their relationship. However
the entanglement is partly concerned with our perception of concepts, not just
their ontology; it starts, after all, from phenomenology
\autocite[21--22]{noeEntanglement2023}. While we view tools as extending our
habitual selves, it seems that we do not view art in the same way. Art seems to
reorganize, rather than extend, our habitual ways of being
\autocite[21--22]{noeEntanglement2023}.

%For instance, consider the difference between how life-changing tools affect us
%and how life-changing art affects us. In the former case, tasks are made easier
%or more efficient, but this does not have to be true for the latter.

Moreover, items on a spectrum may contextualize each other, but they do not have
to change each other. However, an entanglement requires mutual development.
Crucially, insofar as dance is something beyond mere expressive movement, it is
entangled. Because second-order choreography actively loops down to reorganize
our ``spontaneous'' practices, the two cannot be isolated. There is thus no
history of dance outside of its entanglement with Dance
\autocite[32]{noeEntanglement2023}. Yet, if our first-order practices are always
already informed by the second-order, we are left with some critical questions:
what was the first-order all along? Are first-order selves, practices, or
concepts even possible?

It is telling that Noë largely moves away from the language of first- and
second-orders as \textit{The Entanglement} progresses. These terms are
illustrative, but they simplify and distort the entanglement in two ways: first,
the terms ``first-order'' and ``second-order'' imply a chronology. They suggest
a story in which some second-order act of reflection interrupts a pristine,
natural, yet unrefined practice or way of being. However, part of Noë's point is
that no such state exists. Second, and this related to my former concern,
``first'' and ``second'' are absolute terms. In the sense described thus far,
belonging to the first-order could suggest a kind of origin. Despite this, the
entanglement is incompatible with any sort of stable and historical essence
\autocite[32]{noeEntanglement2023}.

To be sure, the language of first- and second-orders is flawed. Regardless, the
source of these flaws is a notable, albeit implicit, feature of the
entanglement: recursivity. As mentioned above, an essential part of going
second-order is putting the first-order on display; in doing so, it ``enables
us to do it all differently'' \autocite[42]{noeEntanglement2023}. This was
apparent in the case of dance and Dance. My contribution, here, is to point out
that Dance's status as ``second-order'' is a relative, rather than absolute or
inherent, property; Dance is only ``second-order'' in relation to dance.
However, nothing about Dance per se is second-order.

For example, consider postmodern art. Within a given artistic practice, one is
likely to encounter postmodernist variations and traditions. Take, for instance,
the example of postmodern choreography. In \textit{Postmodern Dance}, dance
historian Sally Banes argues that postmodern choreography is defined partly by
its relationship with other forms of dance (or, in Noë's terms, Dance). The
relationship she describes is critical and dialectical, with postmodern dance
engaging in an ``anti-modernist, interdisciplinary mingling'' of ``previously
separate spheres'' \autocite[154--155]{banesPostmodern1997}. Notably, as a term,
``postmodern dance'' was first used, not to identify a particular style, but
instead to identify a movement that explicitly rejected modernist styles of
dance \autocite[151]{banesPostmodern1997}.
% I am considering touching on style.

Within Dance, postmodern choreography reacts to other styles of Dance and puts
them on display; it calls on and incorporates other traditions and disciplines,
changing itself in the process. For example, Banes points out that the
politicization of various identities led postmodern dance to foreground and call
on various traditions to create a ``conscious articulation of difference'':

\begin{quote}
    To be sure, some issues of political identity have been and continue to be
    explored and celebrated in dance in a more traditional way ... But in the
    nineties it became clear that many dancers identified themselves as
    bicultural in a specifically avant-garde mode: black and postmodern; Latino
    and postmodern; Asian American and postmodern; gay and postmodern.
    \autocite[154]{banesPostmodern1997}
\end{quote}

Postmodern Dance is informed by the rest of Dance (as well as culture
more broadly), similar to how dance influences choreography as a whole.  Thus,
postmodern choreography goes second-order within Dance, a practice which itself
belongs to the ``second-order.''  This is not to say that modernist, or any
other other styles of Dance, are somehow unthinking or lesser; at a minimum,
modernist choreography and lowercase ``d'' dance are entangled. Moreover, as
previously discussed, ``first-order'' concepts are not naïve. Rather, the point
of this example is to demonstrate that entanglements can be recursive or nested;
the entanglement of dance and Dance was just the beginning.

% I forgot about my attempt to actually define what the first order is...
From the arguments above, I draw two conclusions: first, instead of restricting
ourselves to first- and second-orders, we ought to think of entanglements as
involving an order $n$ and an order $n+1$. It is difficult to pin down the
ultimate origin of any entanglement, if such a thing exists at all. Hence,
labeling something as first-order is a kind of ``second-order'' fallacy where
one takes the substance provided by the first-order for granted without
realizing their position in any other entanglements. In truth, what appears to
be an order $n+1$ practice already belongs to the order $n$ of another
entanglement. For the sake of convenience and legibility, I will use the terms
``lower-order'' and ``higher-order'' instead of ``order $n$'' and ``order
$n+1$,'' respectively. While these terms are less precise, they capture the
concepts' relativity and help separate my interpretation from the simpler first-
and second-order presentation of the entanglement.

Second, the members or objects of an entanglement (such as dance and Dance) are
not ``atomic'' concepts; Entangled concepts are internally complex partly
because they themselves contain entanglements. Consequently, entangled concepts
operate more like spaces with their own internal dynamics and less like static
definitions from a dictionary. This structure reveals the kinds of concepts that
can be entangled; it is not limited to individual practices or objects, but
could extend to entire domains or realms of discourse. Further support for this
interpretation comes from Noë's application of the entanglement; for instance,
he argues that science and philosophy are entangled
\autocite[205]{noeEntanglement2023}.

For Noë, science and philosophy are ultimately one epistemic project. Further,
Noë argues that, in its revolutionary moments, scientists engage in philosophy;
rather than working within a pre-existing framework or developing models,
revolutionary science requires one to change their understanding of a field or
area of research as a whole. For Noë, this kind of work is inherently
philosophical. Moreover, he notes that philosophy needs science as well
\autocite[205]{noeEntanglement2023}. Thus, the entanglement of philosophy and
science, on Noë's view, involves undermining an initial understanding of
philosophy and science as distinct disciplines. The co-dependence of philosophy
and science, as well as their overlapping goals, unites them into one entangled
object. Nonetheless, the conceptual distinctions between these disciplines are
still intact; hence, they are entangled insofar as they are separate and unified
\autocite[205]{noeEntanglement2023}.

% Instead of the paragraph above, I should use the transition I have given
% myself. I should talk more about how Noë views the entanglement between
% science and philosophy. Then, I can contrast it with a more Arendtian view.
However, as I have previously argued, entanglements are nested; as a result,
there is room to zoom out of this particular entanglement. I contend that, from
this new vantage point, one begins to notice that all of the aforementioned
fields are entangled with a broader sphere of common human affairs. This
entanglement is less like that of philosophy and science, and more like that of
dance and Dance. In the latter case, a kind of ``looping down'' characterizes
the entanglement. The ``looping concept'' is an idea Noë borrows from Ian
Hacking, and it is key to understanding this entanglement. In essence, a looping
concept provides a template, identity, or ideal. These ideals then provide
normative standards for what does and does not count as belonging to the
concept. The concept thus ``loops down'' by providing ``self-understanding for
those who fall under them'' \autocite[153]{noeEntanglement2023}.

For example, one's notion of what constitutes good dancing, even in an everyday
sense, might include virtues like grace or expressiveness. Yet, the generation
of such standards requires a kind of higher-order reflection on the habitual act
of dancing. In this case, Dance is looping down to dance. This is the kind of
entanglement I have in mind when I claim that philosophy, science, and even art
are entangled with a broader sphere of human affairs. But what is this sphere,
and what is its entanglement like? To answer these questions, I turn to
political theorist Hannah Arendt. Arendt views common, social human affairs as
political insofar as they are free from the demands of animal necessity, and she
locates this conception of politics in the public sphere
\autocite[157]{arendtPromise2009}.

% TODO: Next I need to develop the exclusive as higher order, explicitly calling
% on Byung-Chul Han. Only after developing the public sphere though.
The public sphere, like the lower-order activity of dance, is close to our
immediate experiences. Consequently, it bears many assumptions that might make
it appear naïve. One notable assumption is that of freedom. Arendt notes that
one encounters freedom, not as an inner quality of their will, but as a result
of their experience in a shared, tangible, outer world. We learn, for instance,
that those bound to necessity or coercion are unfree. One is free, by contrast,
when they assert themselves in the world through ``word and deed''
\autocite[148]{arendtFuture1968}. Crucially, this knowledge is partly implicit;
we do not consciously deliberate to arrive at it. Instead, we recognize it as a
precondition to the political sphere; we hold freedom as a societal standard
without concerning ourselves with its metaphysics
\autocite[152]{arendtFuture1968}.

Nevertheless, according to Arendt, the inner, metaphysical sense of freedom has
its origin in the public sphere. She argues that the problem of free will begins
with the transference of the public concept of freedom to an ``inward domain.''
Once freedom became a topic of philosophical reflection by way of an inward
domain, it lost its public meaning \autocite[146]{arendtFuture1968}. There is an
analogy here between lower-order activities such as dance and the public sphere;
just as dance helps substantiate Dance, the public sphere helps substantiate
philosophy (and science) by providing and originating certain concepts (freedom,
in this case). Here, it is worth noting the epistemic aspects of this emerging
entanglement; the public sphere foregrounds the aforementioned ``word and
deed,'' while science and philosophy are more overtly concerned with truth.

Arendt herself makes sure to emphasize this point, arguing that this difference
in focus can create tension between the public sphere and more truth-oriented
domains. For instance, Arendt argues that factual truths (a kind of empirical,
contingent fact) are, from the political perspective, domineering or
anti-political \autocite[241]{arendtFuture1968}. Similarly, she argues that
philosophical truth (a kind of rational or ``eternal'' truth) becomes opinion
once it enters the public realm \autocite[238]{arendtFuture1968}. To make
something political, in this sense, is to make it subject to debate. Consider
the political discourse surrounding climate change or vaccination; no matter how
settled the science is, the moment one raises these issues in a political
context, they have pushed the issue into the political realm, which invariably
leads to discourse. Consequently, This move from one space to another is not
merely a move ``from one kind of reasoning to another but from one way of human
existence to another'' \autocite[238]{arendtFuture1968}.

The public sphere's collective, active, and discursive nature can put it at odds
with truth-oriented disciplines. As Han observes, ``truth is exclusive and
selective'' \autocite[40]{hanSwarm2017}. To arrive at the truth, one needs time
and the struggle of negativity; one needs a healthy respect for epistemic
authority, as well. Popular discourse often lacks the bearing required for this
kind of engagement. Arendt herself argues that truth-tellers and truth-oriented
institutions must stand outside of politics; the political realm requires
institutions that are outside of its struggles
\autocite[260--261]{arendtFuture1968}. Notably, Arendt places philosophers,
scientists, historians, and artists outside of the political realm
\autocite[259-260]{arendtFuture1968}. While I find this placement too rigid, it
motivates my notion of an exclusive realm.

Art, history, philosophy, and science all have a higher-order relation to the
public sphere; they are to the public sphere as Dance is to dance. Together they
form an exclusive realm that is oriented towards truth. There may be variations
between how members of the exclusive domain understand truth, but unlike the
public sphere, they nonetheless foreground their conceptualization of it. Note
that exclusivity is often institutional. Even art has obviously institutional
aspects, and these motivate, for instance, institutional definitions of art
\autocite[8]{winnerHow2019}. Institutional definitions of art are particularly
noteworthy here; because philosophers tend to be the ones proposing them, they
represent a fully exclusive understanding of art. Of course, none of these
disciplines are, in truth, wholly exclusive. Yet, that is precisely the point;
the public and exclusive spheres are entangled.

For example, artists are a point of entanglement between the public and
exclusive. Arendt points out that artists, like philosophers and scientists,
stand outside the political realm; being an artist is a mode of ``being alone''
\autocite[259--260]{arendtFuture1968}. Artists must turn their attention away
from the world for the process of creation. Moreover, as Noë points out,
detachment and disinterestedness are key features of the aesthetic attitude
\autocite[188--189]{noeEntanglement2023}. Further, it is apparent that many of
our artistic categories have institutional or academic origins. Regardless, this
is only one part of their entanglement. Arendt notes that being an artist is
simply one mode of being \autocite[259--260]{arendtFuture1968}. Moreover, Noë
has already demonstrated how art makes use of lower orders, as in the case of
dance and Dance.

Thus, insofar as one is an artist, they live in an exclusive domain.
Nonetheless, they call on what the public realm originated. In return, they
enrich the public sphere. Crucially, most are never just artists; they
eventually return to the public sphere and embody a common mode of being.
Notably, Nietzsche has argued that, even when an artist studies history, they do
so with an ``eye towards life'' \autocite[100--101]{nietzscheUnfashionable1995}.
In essence, artists, at some level, recognize their debt to the political realm
and pay it gladly. Artists are careful with public concepts and remember where
they came from, so we are receptive to the artist's work; it ``loops down'' and
changes us.

One might characterize this relationship in terms of respect and authority;
lower-order domains grant higher-order domains their authority insofar as the
latter respects what it pulls from the former. In these terms, a lack of respect
shows up as denying the public sphere's centrality. We see glimpses of this
relationship in Arendt's analysis of the problem of free will. Free will, she
argues, was a concept belonging to politics and interpersonal affairs
\autocite[145]{arendtFuture1968}. As mentioned earlier, One encounters freedom
in the public realm.

% NOTE: I do not love the determinism example. I think compatibilsim is more
% popular, and it also does kind of illustrate the point; compatibilsm basically
% says that the public and exclusive concepts of freedom don't interact. Do I
% need a source for this? I think these terms should be widely known enough
% within my audience? Maybe I should find one later anyways. TODO: Find a source
% for compatibilism and determinism claims.
By taking this concept of freedom from the public sphere and placing it into an
``inward domain,'' it becomes subject to ``self-inspection.'' Self-introspection
is an isolated activity, and in it, ``philosophic and metaphysical questions
arise'' \autocite[145]{arendtFuture1968}. Hence, the problem of free will is no
problem of freedom itself; instead, it is a problem that arises once one moves
the concept from the public sphere to the exclusive domain (philosophy, in
particular). For instance, determinism, despite its force in philosophic or
scientific contexts, has not yet looped down in the same way various exclusive
artistic practices have. Moreover, the appeal of compatibilism arguably lies in
its concession to the public sphere; materially, we are determined, but this
does not exclude the possibility of freedom, as freedom operates in a separate
domain.

Even assuming a deterministic view, one can think they have changed their
actions or thoughts, or that they have been changed by some experience, and
these supposed changes can affect politics; they can be the alleged reasons for
action. In contrast, most still cannot help but take themselves and others to be
free, especially regarding legal and moral issues. The problem of free will,
then, is a problem of respect; the public concept of freedom entered the
exclusive domain, but never honored the original concept. Because of this broken
entanglement, deterministic views of free will remain politically impotent.
Here, I argue that certain approaches in cognitive science risk a similar
disconnect.

In \textit{Mental Imagery: Philosophy, Psychology, Neuroscience}, Bence Nanay
offers an interdisciplinary account of mental imagery. Additionally, he argues
that mental imagery founds many human faculties, including memory and the senses
\autocite[viii]{nanayMental2023}. Nanay's arguments exemplify some of the
concerns presented above. First, Nanay claims that {``}`Mental imagery' is a
technical term'' \autocite[3]{nanayMental2023}. He repeatedly emphasizes this
point, later stating that his ``aim here is not to capture what the general
public means by mental imagery'' \autocite[22]{nanayMental2023}. As a result,
Nanay does not think there are ``pre-theoretical considerations that should
persuade us to use the term one way or another''; instead we ought to use the
term in ways that are ``theoretically fruitful'' \autocite[22]{nanayMental2023}.

% NOTE: Two problems with this first view: general (never a clean distinction
% here) and cog sci specific (cognition is common and entangled). Also, I really
% need to mention that part of Nanay's goal is to further our understanding of
% the mind itself.
In terms of the framework developed thus far, Nanay is claiming that mental
imagery is from the exclusive domain; it is not a public concept. If this is
true, then the concept and those who study it are not indebted to the public
sphere. Hence, it is reasonable to prioritize theoretical fruitfulness above all
else. However, the entanglement of the public and exclusive realms means that
this separation is not so clean, especially in the case of cognitive science.
Scientists, as Arendt has pointed out, are humans, and as such eventually return
to the public sphere \autocite[268--269]{arendtFuture1968}. Humans, insofar as
they are members of a common realm, initiate the process of science; this common
realm is where we find humanistic (but non-scientific) motivators such as one's
values, curiosity, or the meaning of their life.

This view aligns with Noë, who claims that the motivations and tools of science
are not themselves scientific \autocite[209]{noeEntanglement2023}. In fact, this
must be the case; science cannot support itself. As Arendt points out, it is an
ideal of modern science to exclude ``all such anthropocentric, that is, truly
humanistic, concerns'' \autocite[266]{arendtFuture1968}. In my view, this ideal
is what motivates another Arendtian concept, one I call functionalization.
Arendt, in her description of liberal and conservative views of history and
politics, notes that it was common at the time to view communism as a new
atheistic ``religion'' \autocite[101--102]{arendtFuture1968}. It was understood
as a religion because, ``socially, psychologically, and `emotionally,'{''} it
performed the same function as religion \autocite[101--102]{arendtFuture1968}.
Arendt also suggests that this same issue exists for authority and violence;
because they can function similarly, they were viewed as identical
\autocite[103]{arendtFuture1968}.

In identifying functionality and ontology, we lose our ability to make
distinctions between concepts, thus muddying our discourse
\autocite[101-102]{arendtFuture1968}. From an Arendtian perspective,
prioritizing theoretical fruitfulness over pre-theoretical meaning risks
functionalizing the concept of mental imagery. While, the term ``mental
imagery'' has a technical origin, the phenomenon underlying mental imagery is
present in the public sphere; it is present, for instance, when one imagines an
object \autocite[3]{nanayMental2023}. It is this qualitative experience and our
use of it that lends the concept its salience. However, by prioritizing
theoretical fruitfulness, Nanay's technical definition separates mental imagery
from this experience.

This separation is often explicit and intentional. For example, Nanay claims
that cutting ties with conscious phenomenology is what made mental imagery a
``respectable concept to study in the empirical sciences of the mind''
\autocite[13]{nanayMental2023}. This has two effects that are particularly
relevant to this discussion. First, as Nanay points out, it makes the concept
quite broad, especially compared to public conceptualizations
\autocite[251]{nanayMental2023}. In line with Arendt's worries about
functionalization, it seems that this treatment of mental imagery does threaten
certain conceptual distinctions \autocite[101--102]{arendtFuture1968}. Second,
it removes agency and conscious experience from the concept \autocite[17,
23]{nanayMental2023}. This latter point relates to Arendt's notion of an
Archimedean point; as science leaves humanistic concerns behind for its own
theoretical development, it takes a perspective on humanity that reduces it to
large-scale biological processes \autocite[279--280]{arendtFuture1968}.

An essential aspect of arriving at this Archimedean point, I would argue, is
functionalization. Functionalization is a way to weaken the relationship between
exclusive domains and the public sphere. While it is by no means a tool of the
exclusive alone, it still offers a means of devaluing experience and agency.
However, as mentioned earlier, for a concept to ``loop down'' and affect public
concepts and being, it must maintain some connection to the political realm. If
cognitive science develops its concepts in spite of discontinuities with their
public forms, it sacrifices the capacity to meaningfully inform our
self-understanding. In the case of art and the cognitive science of art, this
indifference to the public sphere is particularly salient.

When Nanay applies his technical concept of mental imagery to aesthetics, he
frames the observer's engagement with art as a passive event. For example, when
discussing how film directors like Jean-Luc Godard or Luis Buñuel use occlusion
and spatial ambiguity, Nanay argues that these aesthetic choices ``crossmodally
trigger'' indeterminate mental imagery in the viewer
\autocite[242--243]{nanayMental2023}. In his analysis, the aesthetic experience
is described in terms of involuntary, subpersonal processes; the audience's
mental imagery is automatically elicited by the sensory gaps provided by the
artwork. In this framework, the observer is not an active participant engaged in
a discursive cultural practice; instead, they are the coincidental and passive
site of their own perceptual mechanisms.

Alva Noë argues that views like this rely on a ``trigger conception'' of art
\autocite[146]{noeEntanglement2023}. The trigger conception views artworks as
special because they can cause aesthetic experiences in consumers
\autocite[146]{noeEntanglement2023}. However, Noë argues that the trigger
conception relies on a faulty notion of ``aesthetic experiences''; they assume
that aesthetic experiences are momentary episodes in consciousness ``like a
ringing in one's ears'' \autocite[146]{noeEntanglement2023}. For many exclusive
domains, cognitive science in particular, this conception is helpful; if
aesthetic experiences are moments anchored to artistic stimuli, then it is
easier to bring them into a lab and study them. Of course, art's power to
reorganize us weakens this understanding of aesthetic experiences; when we are
changed by art, that change stays even after we stop directly engaging with the
work.

Like functionalization, therefore, the trigger conception of art is a means of
severing a concept's connection to the public sphere. However, unlike
functionalization, the trigger conception of art does not necessarily remove or
downplay parts of the original concept. Instead, it declares that aesthetic
engagement is strictly modal. On this view, some artistic stimuli triggers our
perceptual systems, causing us to enter an aesthetic mode of cognition. Once the
perceptual input has escaped our attention, we leave this mode, returning to a
default state. Under the trigger conception, it is as if aesthetic experiences
are an altered state of consciousness.

% --- Works Cited ---
\newpage
\printbibliography

\end{document}
